\begin{longtable}{
    |>{\RaggedRight\arraybackslash}p{0.319\textwidth}
    |p{0.637\textwidth}|
    } % 0.956
    \caption{Определения цикла в официалльных публикациях международных организаций и государственных органов}
    \label{tab:cycle-definitions} \\
    \LongTableHeader{
        \hline
        \rowcolor{black}
        \textcolor{white}{\textbf{Источник}} &
        \textcolor{white}{\textbf{Определение цикла}} \\
        \hline
    }

    \longcaptioneachpage{Источник: составлено автором}

    NBER (2025). Business Cycle Dating Procedure: Frequently Asked Questions\cite{NBER-2025}
    &
    Деловой цикл понимается как рекуррентные флуктуации совокупной активности между пиками и доньями \\ \hline

    World Bank (2014). Cyclical Drivers of Business Cycles in Emerging Market Economies. Policy Research Working Paper\cite{WB-2014}
    &
    Business cycle – fluctuations in production, trade, and general economic activity over a span of months or years … The term "cycle" is imprecise, as fluctuations are neither regular nor predictable. Expansions and contractions can be measured relative to the long-term GDP growth trend \\ \hline

    IMF (1998). World Economic Outlook, May 1998. Chapter III \textquotedbl{}The Business Cycle, International Linkages, and Exchange Rates\textquotedbl{}\cite{IMF-1998}
    &
    Циклические флуктуации экономической активности рассматриваются как характерная особенность поведения большинства экономик; анализируются как отклонения от тренда ВВП и других макропоказателей \\ \hline

    ECB (2023). The Euro Area Business Cycle and Its Drivers. Occasional Paper No. 354\cite{ECB-2023}
    &
    Определяются классический цикл (чередование периодов рецессии и расширения в уровнях ВВП) и цикл роста (отклонения ВВП от тренда); выделяются четыре фазы: расширение, замедление, сжатие, оживление \\ \hline

    OECD (2011). Business Cycle Measurement and Analysis, Main Economic Indicators methodology\cite{OECD-2011}
    &
    Различаются классический цикл (подход NBER – пики и дна в уровнях активности) и цикл роста (growth cycle — отклонения от тренда); деловые циклы описываются как рекуррентные флуктуации совокупной экономической активности в рыночно-ориентированных экономиках \\ \hline

    Eurostat (2017). Statistics Explained. Glossary: \textquotedbl{}Business cycle\textquotedbl{}\cite{Eurostat-2017}
    &
    A business cycle describes the expansions and contractions of economic activity in an economy over a period of time … expressed as a percentage change of an economic indicator: production, hours worked, employment, GDP, etc. \\ \hline

    Federal Reserve Bank of St. Louis (2023). Page One Economics: All About the Business Cycle: Where Do Recessions Come From?\cite{FRED-2023}
    &
    Business cycle: The fluctuating levels of economic activity in an economy over a period of time measured from the beginning of one recession to the beginning of the next \\ \hline

    Office for Budget Responsibility (2017). Economic Cycles and the Long-term Projections. Technical Paper\cite{OBR-2017}
    &
    Экономические циклы описываются как симметричные или асимметричные отклонения ВВП от долгосрочного тренда; моделируются как чередование периодов, когда ВВП находится выше или ниже тренда, с амплитудой порядка 2\% и длиной волны около 3 лет \\ \hline
\end{longtable}